% TOP_PROBLEM: {"id": 2703, "title": "Kirby Problem 1.44", "category_id": 7, "category": {"id": 7, "name": "topology", "display_name": "Topology", "description": "Properties preserved under continuous deformations.", "slug": "topology", "order_index": 7, "created_at": "2026-07-31T15:26:25.671Z"}, "status": "open"}
% FIRST_POSTED: null
% ENTRY_KIND: statement_only
% Imported from: batch10.tex; UnsolvedMath ID 2703
\documentclass[11pt]{article}
\usepackage[margin=25mm]{geometry}
\usepackage{fontspec}\setmainfont{Latin Modern Roman}
\usepackage{amsmath,amssymb,mathtools}
\usepackage{unicode-math}\setmathfont{Latin Modern Math}
\usepackage{xurl,hyperref}
\hypersetup{colorlinks=true,urlcolor=blue}
\setcounter{secnumdepth}{0}
\setlength{\parindent}{0pt}\setlength{\parskip}{5pt}
\setlength{\emergencystretch}{3em}
\newcommand{\sref}[2]{\hyperlink{source:#1:#2}{[S#2]}}
\newcommand{\cataloguescope}[1]{\paragraph{Recorded target.}#1}
\begin{document}
% UnsolvedMath ID 2703; source catalogue code KP-1.44
\setcounter{section}{2702}
\section{Kirby Problem 1.44}\label{2703}
{\small\sffamily UnsolvedMath ID 2703 · Topology}
\subsection{Definitions and mathematical statement}

\paragraph{Shared notation.}$\mathbb N=\{1,2,\ldots\}$, and $\mathbb Z,\mathbb Q,\mathbb R,\mathbb C$ denote the integers, rationals, reals and complex numbers. Any other conventions are specified locally.


The smooth concordance group $\mathcal C$ consists of oriented knots modulo smooth annuli in $S^3\times[0,1]$, with connected sum as addition and unknot $U$ as zero. A knot is smoothly slice if it bounds a smooth disk in $B^4$, and topologically slice if it bounds a locally flat disk. For a pattern $P\subset S^1\times D^2$, the satellite $P(K)$ uses the preferred longitude; its winding number is the signed homology class in the solid torus. Satellite operations induce functions on $\mathcal C$, not necessarily homomorphisms. For a group $\pi$, put $\pi^{(0)}=\pi$ and $\pi^{(r+1)}=[\pi^{(r)},\pi^{(r)}]$. If $W$ is a compact spin $4$-manifold, let $W^{(r)}$ be the cover corresponding to $\pi_1(W)^{(r)}$. On $H_2(W^{(r)};\mathbb Z)$ the equivariant intersection and self-intersection forms, denoted $\lambda_r,\mu_r$, record intersections with their deck-transformation labels in $\mathbb Z[\pi/\pi^{(r)}]$ (with the usual self-intersection quotient). An $r$-Lagrangian is a submodule on which both forms vanish and whose image in $H_2(W;\mathbb Z)$ is a half-rank direct-summand Lagrangian for the ordinary intersection form. It has $r$-duals if it has generators represented by $r$-surfaces $\ell_i$ and there are $r$-surfaces $d_j$ with $\lambda_r(\ell_i,d_j)=\delta_{ij}$ and $\operatorname{rank}H_2(W)=2\#\{\ell_i\}$. Here an $r$-surface has fundamental group mapping into $\pi^{(r)}$, and its chosen lift represents the specified class. A knot is $r$-solvable if its zero-surgery bounds such a spin $W$ inducing an isomorphism on $H_1$ and admitting an $r$-Lagrangian with $r$-duals. It is $(r+\tfrac12)$-solvable if $W$ admits an $(r+1)$-Lagrangian whose projection has $r$-duals. Write $\mathcal F_n$ for these concordance classes, $n\in\tfrac12\mathbb N_0$. For an integer $r\geq0$, a knot is $r$-positive if it bounds a null-homologous smooth disk $\Delta$ in a compact oriented simply connected smooth $V$ with $\partial V=S^3$, positive definite intersection form, and a basis of $H_2(V)$ represented by disjoint closed oriented surfaces $S_i\subset V\setminus\Delta$ with $\pi_1(S_i)$ mapping into $\pi_1(V\setminus\Delta)^{(r)}$. Define $r$-negative by a negative definite form. Let $\mathcal T_r$ be the subgroup of topologically slice classes that are both $r$-positive and $r$-negative, allowing different $V,\Delta$ in the two conditions. These are the solvable and bipolar filtrations, with definitions from the original construction papers.
\paragraph{Statement recorded in the supplied source.}
(a) If $[K]\in\bigcap_{n\in\frac12\mathbb N_0}\mathcal F_n$, must $K$ be topologically slice? (b) If $[K]\in\bigcap_{r\in\mathbb N_0}\mathcal T_r$, must $K$ be smoothly slice? \sref{2703}{1}\sref{2703}{2}\sref{2703}{3}\sref{2703}{4}
\subsection{Short English statement}
Ask whether satisfying every solvable-filtration condition forces topological sliceness, and whether every bipolar condition forces smooth sliceness.
\subsection{Sources}
\begin{description}
\item[\hypertarget{source:2703:1}{S1}] ULAM.AI and the UnsolvedMath Contributors, \emph{UnsolvedMath: A Curated Collection of Open Mathematics Problems}, record 2703 (KP-1.44). CC BY 4.0. \url{https://huggingface.co/datasets/ulamai/UnsolvedMath}.
\item[\hypertarget{source:2703:2}{S2}] R. \.{I}nan\c{c} Baykur, Robion C. Kirby and Daniel Ruberman, K3 -- A New Problem List in Low-Dimensional Topology, author preliminary edition (2026), Problem 1.44 and its remarks. \url{https://math.berkeley.edu/sites/default/files/surv-295-ruberman-watermarked-author-pdf.pdf}.
\item[\hypertarget{source:2703:3}{S3}] Tim D. Cochran, Kent E. Orr and Peter Teichner, \emph{Knot concordance, Whitney towers and L2-signatures, Annals of Mathematics 157 (2003), 433--519}, Definitions 1.2, 8.1, 8.3, 8.5 and 8.7; Theorem 8.8. \url{https://arxiv.org/abs/math/9908117v2}.
\item[\hypertarget{source:2703:4}{S4}] Tim D. Cochran, Shelly Harvey and Peter Horn, \emph{Filtering smooth concordance classes of topologically slice knots, Geometry and Topology 17 (2013), 2103--2162}, Definitions 2.1, 2.2, 2.4 and 2.6. \url{https://arxiv.org/abs/1201.6283}.
\end{description}
\label{end:2703}
\end{document}
